\begin{abstract}

A simulação numérica é ferramenta central de análise e projeto em Engenharia Mecânica, mas seu uso no ensino esbarra em duas dificuldades recorrentes: as ferramentas mais capazes exigem licenciamento, infraestrutura e configuração de ambiente, e raramente expõem de forma transparente a formulação matemática e o código de cada simulação. Este trabalho se propõe a oferecer um ambiente de simulação didático que reduza essa barreira de acesso, sem abrir mão do rigor técnico nem da transparência da implementação. Apresenta-se o \textit{MinervaMesh}, plataforma web de código aberto que resolve Equações Diferenciais Parciais de interesse em Engenharia Mecânica pelo Método dos Elementos Finitos (MEF) e pelo Método das Diferenças Finitas (MDF). O objetivo é a plataforma em si: a entrega das simulações integralmente no navegador, em arquitetura \textit{client-side} sobre \textit{WebAssembly} e \textit{PyScript}, sem instalação e com código-fonte Python inspecionável. Como fundamento técnico, os métodos numéricos são implementados e verificados contra soluções analíticas e, nos casos sem solução fechada, contra \textit{benchmarks} consagrados da literatura. Os resultados mostram que as equações discretizadas reproduzem as soluções de referência e que a plataforma é adequada ao ensino, à prototipagem e a demonstrações de mecânica computacional.

\end{abstract}
